\chapter{Representation-separation diagnostic для diagonal connectors}

Universal-singlet no-go определяет требуемую характеристику следующего
объекта: он обязан различать required Goldstone \(G\) и unwanted mode
\(X\). В опубликованном разложении они имеют квантовые числа
\begin{equation}
 G:\quad (I_R,R_4)=(1,\mathbf{10}),
 \qquad
 X:\quad (I_R,R_4)=(0,\mathbf6).
\end{equation}

\section{Casimir sensitivity test}

Для \(SU(2)_R\)
\begin{equation}
 C_2(1)=2,
 \qquad C_2(0)=0,
\end{equation}
а для двухиндексных представлений \(SU(4)\)
\begin{equation}
 C_2(\mathbf{10})=\frac92,
 \qquad
 C_2(\mathbf6)=\frac52.
\end{equation}
Поэтому минимальный Casimir diagnostic даёт
\begin{center}
\begin{tabular}{lccc}
\toprule
operator & eigenvalue on \(G\) & eigenvalue on \(X\) & \(X-G\) \\
\midrule
identity & \(1\) & \(1\) & \(0\) \\
\(C_2[SU(2)_R]\) & \(2\) & \(0\) & \(-2\) \\
\(C_2[SU(4)]\) & \(9/2\) & \(5/2\) & \(-2\) \\
sum of Casimirs & \(13/2\) & \(5/2\) & \(-4\) \\
\bottomrule
\end{tabular}
\end{center}

Следовательно, representations различимы group-theoretically: identity их
не различает, а quadratic Casimirs различают. Это не является mass matrix
конкретного adjoint VEV. Для фиксированного breaking direction массы зависят
от weights, tensor contractions и coefficients полного potential, а не
автоматически от полного \(C_2(R)\).

\section{\texorpdfstring{Внешний кандидат: единый \(\mathbf1\oplus\mathbf{15}\) block}{Внешний кандидат: единый 1+15 block}}

Наиболее полезен не формально наименьший \(SU(2)_R\) adjoint, а Hermitian
four-color diagonal block
\begin{equation}
 S_4=s_0\mathbf1_4+s_A T^A,
 \qquad
 S_4\in\mathbf1\oplus\mathbf{15}.
\end{equation}
Причина состоит в объединении двух задач:
\begin{itemize}
 \item trace component \(s_0\) может быть общей scale/dilaton mode;
 \item \(\mathbf{10}\) и \(\mathbf6\) имеют разные Casimirs, поэтому
 representation-sensitive contractions не запрещены group theory;
 \item такой weak-singlet block отсутствует в general fundamental branch,
 но присутствует в constrained composite branch.
\end{itemize}

\begin{proposition}[Первый target-кандидат]
Four-color block \(\mathbf1\oplus\mathbf{15}\) является первым найденным
target-объектом, который может нести общий transmutation scale и проходит
representation-separation diagnostic. Реальное mass splitting и наличие
такого block в project finite geometry не доказаны.
\end{proposition}

Это ещё не vacuum theorem. В general fundamental branch поле \(\Sigma\)
несёт \((2_R,2_L,1+15_4)\). В composite branch отдельный
\((1_R,1_L,15_4)\) появляется после first-order restriction, но его нельзя
считать автоматически выведенным spontaneous remnant. Следующий blind gate
должен воспроизвести обе ветви из одной matrix convention и затем
проверить:
\begin{enumerate}
 \item существование stationary adjoint VEV;
 \item положительность Hessian вне точных Goldstone directions;
 \item отсутствие лишних massless colored states;
 \item отсутствие недопустимых high-scale fermion masses;
 \item Coleman--Weinberg coefficient \(B>0\) для trace mode;
 \item полученный \(m_R\) до gauge-running comparison.
\end{enumerate}

\begin{nogocriterion}[Adjoint coefficient fit]
Ненулевая Casimir difference является только диагностиком различимости.
Коэффициенты identity и adjoint invariants запрещено выбирать из желаемого
scalar spectrum. Они должны следовать из trace normalization одного
finite-Dirac block.
\end{nogocriterion}

Машинный menu сохранён в
\path{s2t_v4_pati_salam_diagonal_connector_menu.py} и
\path{s2t_v4_pati_salam_diagonal_connector_menu_results.json}.